\documentclass[11pt,a4paper]{article}
% Compile with XeLaTeX (fontspec and xeCJK).
\usepackage{fontspec}
\setmainfont{LibertinusSerif}[Extension=.otf,UprightFont=*-Regular,BoldFont=*-Bold,ItalicFont=*-Italic,BoldItalicFont=*-BoldItalic]
\setsansfont{LibertinusSans}[Extension=.otf,UprightFont=*-Regular,BoldFont=*-Bold,ItalicFont=*-Italic]
\usepackage{xeCJK}
\setCJKmainfont{Hiragino Mincho ProN}
\xeCJKDeclareCharClass{Default}{"00B7, "2014, "2018 -> "201F, "2026}
\usepackage[english]{babel}
\usepackage[a4paper,margin=2.5cm,top=2.3cm]{geometry}
\usepackage{graphicx}
\usepackage{xcolor}
\usepackage{titlesec}
\usepackage{fancyhdr}
\usepackage[hidelinks]{hyperref}
\definecolor{accent}{RGB}{31,73,125}
\definecolor{muted}{RGB}{93,104,119}
\hypersetup{colorlinks=true,urlcolor=accent,linkcolor=accent,pdfauthor={Romain Fernex},pdftitle={The origins of Korea and the kingdom of Goguryeo}}
\providecommand{\tightlist}{\setlength{\itemsep}{0pt}\setlength{\parskip}{0pt}}
\providecommand{\ul}[1]{\underline{#1}}
\setcounter{secnumdepth}{0}
\titleformat{\section}{\large\bfseries\color{accent}}{}{0em}{}
\setlength{\parindent}{1.2em}
\setlength{\parskip}{0.35em}
\pagestyle{fancy}
\fancyhf{}
\fancyhead[L]{\small\color{muted}Romain Fernex}
\fancyhead[R]{\small\color{muted}\itshape The origins of Korea and the kingdom of Goguryeo}
\fancyfoot[C]{\small\color{muted}\thepage}
\renewcommand{\headrulewidth}{0pt}
\newcommand{\articlefigure}[2]{%
  \begin{center}
  \includegraphics[width=0.55\linewidth,height=6.5cm,keepaspectratio]{#1}%
  \ifx\relax#2\relax\else\par\smallskip{\small\itshape\color{muted}#2\par}\fi
  \end{center}}

\begin{document}
\thispagestyle{empty}
\noindent{\small\sffamily\color{accent}\MakeUppercase{History of Korea · Part III}}\par\medskip
\noindent{\LARGE\bfseries The origins of Korea and the kingdom of Goguryeo\par}
\smallskip\noindent{\large\itshape Between nationalist narratives and historical realities\par}
\medskip
\noindent{\color{accent}\rule{\linewidth}{0.8pt}}\par\smallskip
\noindent{\small\color{muted}Romain Fernex · 2022 · Article in English}\par\bigskip

In the preceding articles, we analysed the beginnings and assertion of Goguryeo as a political entity in its own right on a Korean peninsula nevertheless largely dominated by the influence of successive great Chinese dynasties, which indeed overcame it, at least temporarily. We must now examine how it recovered after the Chinese conquest and, above all, describe in greater detail how it fitted into a renewed geopolitical order whose main dynamics were dictated no longer only by China but by genuine new regional powers, emerging after China's weakening in the second half of the 3rd century. Many historians call this the ``Three Kingdoms'' period, a name largely derived from two foundational texts, although written much later: the Samguk Sagi (Chronicle of the Three Kingdoms) and the Samguk Yusa (Memorabilia of the Three Kingdoms).

\articlefigure{media/media/image1.png}{Map of the Korean peninsula in AD 375}

First, how could Goguryeo, decapitated after its capital fell to Wei China\textbf{{[}1{]}} in the mid-3rd century, rise from its ashes? The main reason for this resurgence of influence lies in the geopolitical context of the time, particularly the power vacuum left by internecine wars between dynasties seeking power in China. Indeed, the Chinese Jin dynasty, which succeeded the Wei dynasty in 266, faced a series of civil wars from 291 onwards, known to historians as the ``War of the Eight Princes''. As its name suggests, this saw several princes of the Jin dynasty fight to seize power in place of Emperor Hui, who had a severe developmental disorder and was therefore considered unfit to rule. Barely five years after the end of the War of the Eight Princes, Jin China collapsed in favour of numerous barbarian kingdoms that divided its territory. This began a dark age for China, corresponding to the Sixteen Kingdoms period and lasting until 439. Thus, from the mid-3rd century, Goguryeo troops that survived the fall of its capital took advantage of China's chaos to recapture East Okjo and then East Ye and re-establish its presence across much of its former territory. This reconquest continued in 303, accompanied by renewed expansionism through the capture of the Lelang and Daifang commanderies, long central to maintaining Chinese authority in the region.

\articlefigure{media/media/image2.jpeg}{Map of China's Sixteen Kingdoms in AD 350}

But Goguryeo's rebirth did not take place alone. It was also at this time\textbf{{[}2{]}} that three other kingdoms emerged, two of which helped give the Three Kingdoms period its name. First came Baekje, born in the Mahan region of the southern peninsula and probably originating among Han populations to which it was culturally close. In the southeast, in the Kyeongju basin, the Kim royal lineage founded Silla, which would emerge victorious at the end of the Three Kingdoms period, although we will return to this later. Finally, modest in size but also relatively important to the regional balance, Kaya occupied the far south of the peninsula in the Naktong River valley. Its population included a significant proportion of Wa\textbf{{[}3{]}}, the founding ethnic group of Japan and particularly of the Yamato kingdom.

The early 4th century thus differed from preceding periods through the emergence of two major new regional powers, Baekje and Silla, and the great return of Goguryeo, finally escaping competition from China. It was during this period that Goguryeo emerged as the region's pre-eminent power, and this new prosperity was accompanied by major economic, cultural and political transformations. We must now analyse these to understand their extent and hence their relevance to the origins of this sense of Koreanness, at the heart of the debate we are examining.

Let us first look at Goguryeo's new political organisation after several decades in China's shadow. Indeed, this period of Chinese ``domination'' saw the end, or at least the weakening, of the system of principalities in favour of a provincial system with provincial officials relaying Chinese central power and, after the reconquest, royal power. The eumnaks, already disappearing, vanished completely, and direct control over the population finally passed to the royal family. The five Bus were thus reorganised into provinces\textbf{{[}4{]}}, each headed by a Yoksal (governor) exercising authority over several towns led by Cheoryogu (magistrates), assisted by Karadal (aldermen). At the bottom of the administrative hierarchy, mainly in rural areas, Rucheo acted as prefects. Finally, all these officials were accountable to the king, the sole holder of authoritative and coercive power.

This substantial centralisation of power enabled the completion of tax reform, replacing the primitive model of tribute demands with systematic and, above all, progressive taxation (to a certain extent). Royal power thus became the sole holder of taxation rights in place of lords and other large landowners. Although it occasionally granted certain high-ranking officials or deserving vassals a sigeup, a plot of land on which the lord could also collect taxes, it received the vast majority of tax revenue. This taxation targeted agricultural income first and was restricted to independent adult men. Faced with particularly significant processes of social diversification and stratification at the time, a genuine consideration of taxpayers' economic circumstances developed. The appearance of social assistance programmes, allowing certain struggling households to receive grain rations in exchange for a debt to the government, thus coincided with this awareness and aimed to provide better care for the kingdom's human capital.

Also domestically, the relocation of the capital to Pyongyang in 343, following repeated attacks by Former Yan China, enabled the expansion of a universal culture specific to Goguryeo despite numerous borrowings from China. This development began with the acceptance of Buddhism as a pillar of this pan-Goguryeo culture, but also with the systematisation of the Chumong myth, which definitively incorporated all Goguryeo monarchs into the kingdom's official and, above all, legendary history. This culture nevertheless had several variants: one can distinguish that developing in the north under Pyongyang's influence from that appearing in the south, still influenced by Goguryeo's former capital in Jian. In many regions, Taoism thus mingled with Buddhism and beliefs indigenous to this part of the peninsula. In this respect, studying the murals in many tombs of Goguryeo rulers and aristocrats, with which we began our discussion, allows us to appreciate the artistic movements that passed through Goguryeo. This evolution is evident both in the materials used to create these sumptuous frescoes and in the iconography dominating each period and region. Thus, during this first period, extending from the early 4th to the mid-6th century, many paintings depict everyday life or the four pioneering deities of Korean mythology. Several elements of Buddhist iconography, such as the lotus flower, also recur in the burial chambers of the period.

\articlefigure{media/media/image3.jpeg}{Mural depicting a hunting scene,}

\emph{found in a tomb dating from the Goguryeo period}

Thus, describing Goguryeo at this time as a mere administrative division of China seems largely unwarranted, or at least particularly misleading.

Nevertheless, now that we have seen how Goguryeo differed from China and other regional powers in many respects, an important instrument of power remains to be analysed: its military strength. Indeed, Goguryeo rapidly emerged as Korea's almost unchallenged regional military power, and above all one native to the peninsula, for nearly a hundred years after centuries of Chinese influence. Let us now see how Goguryeo managed to place itself at the head of the regional geopolitical order.

From 373, under King Sosurim\textbf{{[}5{]}}, and after a sometimes tumultuous beginning with repeated Yan invasions, Goguryeo managed to pacify its relations with Qin China and could finally turn its dreams of expansion southward. Although it had already successfully attacked Silla in 345, obtaining an agreement in 348 in which the latter recognised Goguryeo's military superiority, it now sought to conquer Baekje, which had resisted it until then. This resurgence was marked by the successful capture in 375 of Sugok fortress, held by Baekje, which attempted a retaliatory campaign that failed completely. Yet the true turning point in Goguryeo's expansionist policy came in 391 with the accession of King Kwang-Gae-Toh, Goguryeo's most famous ruler after Chumong. He launched a wave of successful attacks on Baekje and, according to the Samguk Sagi, captured nearly 58 fortresses and 700 villages in less than a decade. At the same time, Goguryeo could count on the unwavering support of Silla, which appeared in many respects to be its vassal. A policy of hostage diplomacy\textbf{{[}6{]}} was also established between the two powers.

\articlefigure{media/media/image4.jpeg}{Portrait of King Kwang-Gae-Toh}

This relationship of vassalage peaked in 400 when, during an attack on Silla by Japanese marauders led by Baekje and its faithful ally Kaya, Goguryeo came to Silla's aid and drove the enemy forces out of its territory. Following this event, Goguryeo troops remained stationed in Silla, and the hostage diplomacy already in place expanded significantly. However, tensions soon marked this relationship, particularly over the garrisons present in Silla, culminating in the accession of a king hostile to Goguryeo. Worse still, Silla drew closer to Baekje, and in the second half of the 5th century the two powers signed an agreement to counter Goguryeo's unstoppable rise.

In 450, an incident at Hasulla fortress, in which a Goguryeo soldier was killed by a Silla commander, revived these tensions, leading four years later to a pre-emptive Goguryeo attack on Silla. Events then followed in rapid succession, and in 464 the massacre of Goguryeo forces based in Silla's capital definitively ended any alliance between the two powers. Eleven years later, Goguryeo captured Baekje's capital, Hanseon, whose power declined with the outbreak of civil war and the assassination of its ruler, while also annexing part of Silla's territory. Goguryeo, now dominating the vast majority of the peninsula, was at the height of its power and negotiated as an equal with Northern Wei China.

But from the end of the 5th century, this dynamic began to reverse. Silla, which had previously intervened relatively little in the region's conflicts, including within its alliance with Baekje, adopted a much more interventionist policy and annexed most of Kaya in 532 to extend its influence. A little later, when Goguryeo attacked Baekje's Toksan fortress, Silla deployed its army and defeated it. Yet it was 551 that marked the definitive end of Goguryeo's golden age, with the capture of the Han River basin (the river flowing through present-day Seoul) by the Silla-Baekje alliance.

We will return to these events and those that followed in the next article, which will examine Goguryeo's final transformation and its fall in favour of Silla. We will also take that opportunity to conclude our analysis of the controversy between China and Korea over Goguryeo's legacy and nature.

Romain Fernex, dual-degree student at ESCP and Korea University

\textbf{Bibliography:}

\begin{itemize}
\item
  Noh Taedon, \emph{Korea's Ancient Koguryo Kingdom : A Socio-political History}, global oriental, 2014
\item
  Jeon Ho-tae, \emph{Goguryeo Culture and Mural Paintings}, The Review of Korean Studies, Vol.7 n°4, 2004
\item
  Kim Chang-seok, \emph{Goguryeo Society and Its Economy}, The Review of Korean Studies, Vol.7 n°4, 2004
\item
  Michael J.Seth,~\emph{A history of Korea -- From Antiquity to the present}, 2010, Rowman \& Littlefield Publishers
\item
  Kyung Moon Hwang,~\emph{A History of Korea}, 3rd edition, 2021, Palgrave
\item
  Kim Taesik\emph{,~Koguryeo and Gaya -- contacts and consequences}, Hungik universitt, journal of Northeast Asian History, Vol. 4 n°1 (p41-82), 2007
\item
  Jung Woon Yong,~\emph{Trends in Koguryo's relationship with Paekche and Silla during the 4th -- 7th centuries}, International Journal of Korean history, Vol. 8, 2006
\end{itemize}

\textbf{{[}1{]}} This refers to China under the Cao Wei dynasty, lasting from AD 220 to 266. It should not be confused with China under the Northern Wei (386-535) or Western Wei (535-557) dynasties, which appeared much later and are also discussed further on in this article.

\textbf{{[}2{]}} Although the Samguk Sagi dates Silla's and Baekje's foundations to 57 BC and 18 BC respectively, no historical evidence currently corroborates this information, and most historians consider these kingdoms' emergence to be well before the date given for Goguryeo in that same work, which is itself relatively plausible. Thus, Silla's and Baekje's emergence was probably contemporary with Goguryeo's rebirth rather than its foundation.

\textbf{{[}3{]}} The Nihon Shoki also mentions a significant presence of these Wa peoples on the peninsula through the mythical kingdom of Minima, said to have belonged to Yamato and to have been lost to Silla. According to historians, however, this was largely imagined, and the existence of a fully fledged Wa kingdom on the peninsula is particularly difficult to prove given current knowledge of this period.

\textbf{{[}4{]}} Several competing theories exist about what exactly these new administrative provinces designated. For example, some suggest Goguryeo was divided between Naepyeong, the area directly influenced by the capital, and Waepyeong, covering the rest of the territory, itself divided into five provinces, each headed by a Yoksal. Others, however, consider that there were ten provinces in total, hence ten Yoksals, and so on\ldots{}

\textbf{{[}5{]}} It was also under this king that Buddhism was officially accepted in Goguryeo and the first legal code was published, laying the foundations of Goguryeo as it existed at its peak.

\textbf{{[}6{]}} Hostage diplomacy involves a country sending high-ranking young people, often members of a royal lineage, to an allied country as guarantees to prevent the alliance from being broken. This same policy would eventually work against Goguryeo by fuelling the division between Silla's pro-Goguryeo and anti-Goguryeo camps.

\end{document}
